%%FILE blueprint/src/macros/common.tex
% Put any macro needed for the project here. Those macros for both the web and print versions of the blueprint.

\newcommand{\C}{\mathbb{C}}
\newcommand{\Z}{\mathbb{Z}}
\newcommand{\N}{\mathbb{N}}
\newcommand{\A}{\mathbb{A}}
\newcommand{\Q}{\mathbb{Q}}
\newcommand{\R}{\mathbb{R}}
\newcommand{\F}{\mathbb{F}}
\newcommand{\dist}{\text{ dist }}
\newcommand{\Qp}{\mathbb{Q}_p}
\newcommand{\Ql}{\mathbb{Q}_\ell}
\newcommand{\Qbar}{\overline{\Q}}
\newcommand{\Qpbar}{\overline{\Q}_p}
\newcommand{\Qlbar}{\overline{\Q}_\ell}
\newcommand{\bbC}{\mathbb{C}}
\newcommand{\GQ}{\Gal(\Qbar/\Q)}
\newcommand{\GQp}{\Gal(\Qpbar/\Qp)}
\newcommand{\GQl}{\Gal(\Qlbar/\Ql)}
\newcommand{\m}{\mathfrak{m}}
\newcommand{\GK}{\Gal(K^{\sep}/K)}
\newcommand{\GN}{\Gal(\overline{N}/N)}
\newcommand{\Kbar}{\overline{K}}
\newcommand{\Qhat}{\widehat{\Q}}
\newcommand{\calO}{\mathcal{O}}
\newcommand{\bbH}{\mathbb{H}}
\newcommand{\p}{{\mathfrak{p}}}
\newcommand{\rhobar}{\overline{\rho}}
\newcommand{\Zhat}{\widehat{\Z}}
\DeclareMathOperator{\Gal}{Gal}
\DeclareMathOperator{\avoid}{avoid}
\DeclareMathOperator{\Aut}{Aut}
\DeclareMathOperator{\GL}{GL}
\DeclareMathOperator{\PGL}{PGL}
\DeclareMathOperator{\PSL}{PSL}
\DeclareMathOperator{\SL}{SL}
\DeclareMathOperator{\Spec}{Spec}
\DeclareMathOperator{\sep}{sep}
\DeclareMathOperator{\ab}{ab}
\DeclareMathOperator{\tr}{tr}
\DeclareMathOperator{\Hom}{Hom}
\DeclareMathOperator{\Frob}{Frob}


\tikzset{
            carc/.style args={#1:#2:#3}{
              insert path={+(#1:#3) arc (#1:#2:#3)}
            }
          }
%%FILE blueprint/src/macros/web.tex
% Those macros are used for the web version of the blueprint.

\newtheorem{theorem}{Theorem}[chapter]
\newtheorem*{theorem*}{Theorem}
\newtheorem{proposition}[theorem]{Proposition}
\newtheorem{lemma}[theorem]{Lemma}
\newtheorem{sublemma}[theorem]{Sub-lemma}
\newtheorem{corollary}[theorem]{Corollary}

\theoremstyle{definition}
\newtheorem{definition}[theorem]{Definition}
\newtheorem{proplem}[theorem]{Proplem}

\theoremstyle{remark}
\newtheorem{remark}[theorem]{Remark}
\newtheorem{example}[theorem]{Example}

\newcommand{\RNum}[1]{#1}
%%FILE blueprint/src/macros/print.tex
% Those macros are used for the print version of the blueprint.

\declaretheorem[numberwithin=chapter]{theorem}
\declaretheorem[numbered=no]{theorem*}
\declaretheorem[sibling=theorem]{proposition}
\declaretheorem[sibling=theorem]{corollary}
\declaretheorem[sibling=theorem]{remark}
\declaretheorem[sibling=theorem]{lemma}
\declaretheorem[sibling=theorem]{definition}
\declaretheorem[sibling=theorem]{problem}
\declaretheorem[sibling=theorem]{example}

% We neutralise the Plastex commands
\newcommand{\uses}[1]{}
\newcommand{\proves}[1]{}
\newcommand{\discussion}[1]{}
\newcommand{\lean}[1]{}
\newcommand{\leanok}{}
\newcommand{\mathlibok}{}
\newcommand{\notready}{}

\newcommand{\RNum}[1]{\uppercase\expandafter{\romannumeral #1\relax}}

%%FILE blueprint/src/BA_Titelseite.sty
\NeedsTeXFormat{LaTeX2e}
\ProvidesPackage{BA_Titelseite}[2010/04/25]

%%%%%%%%%%%%%%%%
% Titelseite:
%%%%%%%%%%%%%%%%
\newcommand*{\betreuer}[1]{\def\@betreuer{#1}}
\betreuer{}
\newcommand*{\zweitgutachter}[1]{\def\@zweitgutachter{#1}}
\zweitgutachter{}
\newcommand*{\ausarbeitungstyp}[1]{\def\@ausarbeitungstyp{#1}}
\ausarbeitungstyp{}
\newcommand*{\geburtsdatum}[1]{\def\@geburtsdatum{#1}}
\geburtsdatum{}
\newcommand*{\geburtsort}[1]{\def\@geburtsort{#1}}
\geburtsort{}
\newcommand*{\institut}[1]{\def\@institut{#1}}
\institut{}

\renewcommand\maketitle{\begin{titlepage}%
  \let\footnotesize\small
  \let\footnoterule\relax
  \let \footnote \thanks
  \null\vfil

\begin{center}%

\parbox{10cm}{\centering\huge\bfseries \@title \par}\\
\vspace{1em}
{\Large
\vspace{1em}
\@author}\\
\vspace{1em}
Geboren am \@geburtsdatum \ in \@geburtsort\\
\vspace{1em}
{\large \@date}
\vspace{10em}

{\large \@ausarbeitungstyp}\\
\vspace{1em}
{\large \@betreuer}\\
\vspace{1em}
{\large \@zweitgutachter}\\
\vspace{1em}
\centerline{{\large\scshape \@institut}}
\vspace{10em}

\centerline{{\large\scshape Mathematisch-Naturwissenschaftliche Fakult\"at der}}
\vspace{1em}
\centerline{{\large\scshape Rheinischen Friedrich-Wilhelms-Universit\"at Bonn}}

\end{center}

\vfil\null
\clearpage
\thispagestyle{empty}\mbox{}
\clearpage
\pagenumbering{arabic}

\end{titlepage}%

  \setcounter{footnote}{0}%
  \global\let\maketitle\relax
  \global\let\@author\@empty
  \global\let\@date\@empty
  \global\let\@title\@empty
  \global\let\title\relax
  \global\let\author\relax
  \global\let\date\relax
  \global\let\and\relax
}
% END Titelseite

\endinput
